\documentclass[11pt,a4paper]{article}
\usepackage[margin=1in]{geometry}
\usepackage{times}
\usepackage{latexsym}
\usepackage{graphicx}
\usepackage{amsmath}
\usepackage{amssymb}
\usepackage{booktabs}
\usepackage{multirow}
\usepackage{xcolor}
\usepackage{enumitem}
\usepackage{tabularx}
\usepackage[colorlinks=true,citecolor=blue,linkcolor=blue]{hyperref}
\usepackage{natbib}
\bibliographystyle{plainnat}

\newcommand{\ours}{CRD-Diag}
\newcommand{\disagree}{\textsc{Disagree}}
\newcommand{\targeted}{\textsc{Targeted}}
\newcommand{\blind}{\textsc{Blind}}

\title{Cross-Representation Disagreement as a Diagnostic Signal\\for LLM Formalization Errors}

\author{Anonymous}

\begin{document}
\maketitle

% ============================================================
\begin{abstract}
Neurosymbolic reasoning pipelines that translate natural language into formal representations suffer from formalization errors as a dominant failure mode. We present the first systematic empirical analysis of \emph{cross-representation disagreement}---the phenomenon where the same LLM produces conflicting answers when formalizing a problem into two distinct representations (first-order logic vs.\ executable code). On three benchmarks spanning logic and mathematics (FOLIO, ProofWriter, MATH Level~5), we show that disagreements occur at a substantial rate (37--42\%) and serve as a high-precision diagnostic signal: 81--87\% of disagreements identify genuine formalization errors. Disagreement patterns enable automatic error localization with 82--85\% accuracy and reveal a striking domain-dependent asymmetry---FOL dominates in logic domains (7--9:1 error ratio) while code dominates in mathematics (2.3:1 inverted). Building on these findings, we develop a lightweight diagnosis-guided repair framework that leverages error type classification to achieve +10.8pp improvement over blind regeneration ($p=0.008$). We further characterize the boundary of this approach: specification-level errors producing correlated failures constitute its ceiling, a finding that deepens understanding of neurosymbolic failure modes.
\end{abstract}

% ============================================================
\section{Introduction}
\label{sec:intro}

Large language models (LLMs) are increasingly used in neurosymbolic reasoning pipelines, where natural language problems are translated into formal representations---first-order logic (FOL), executable Python code, or solver-compatible formats---for downstream verification or execution \cite{pan2023logic,ye2024satlm}. A dominant failure mode in these pipelines is \emph{formalization error}: the LLM's translation introduces subtle corruptions that cause the formal representation to diverge from the intended semantics of the original problem.

Existing approaches address this challenge along three dimensions. \emph{Representation selection} methods choose among multiple formalizations based on confidence or agreement signals \cite{hybridmind2024,adaptive2024}. \emph{Within-representation verification} methods check internal consistency through roundtrip translation or self-verification \cite{vericot2024,verge2024,roundtrip2026}. \emph{Agreement filtering} methods discard solutions where different representations disagree \cite{liang2024cotpot,prove2024}. However, none of these approaches exploit cross-representation disagreement for \emph{error categorization} and \emph{targeted repair}.

We observe that when the same LLM formalizes a reasoning problem into two distinct representations, disagreements are not noise---they are structured diagnostic signals. The pattern of disagreement (which representation errs, how it errs, and on what problem types) carries rich information about the nature of formalization failures.

This paper makes three contributions:

\begin{enumerate}[leftmargin=*,itemsep=2pt]
\item \textbf{Disagreement as high-precision diagnostic signal.} On three benchmarks (FOLIO, ProofWriter, MATH Level~5), cross-representation disagreement occurs at 37--42\% and identifies genuine errors with 81--87\% precision (\S\ref{sec:main-results}).

\item \textbf{Error localization, classification, and domain-dependent asymmetry.} Disagreement patterns enable 82--85\% error localization accuracy and reveal a systematic asymmetry: in logic domains, FOL representations exhibit 7--9$\times$ lower error rates than code; in mathematics, code (SymPy) dominates with a 2.3:1 inverted ratio (\S\ref{sec:asymmetry}).

\item \textbf{Boundary characterization of diagnosis-guided repair.} Multi-round, taxonomy-guided repair achieves +10.8pp over blind regeneration. We characterize the method's ceiling: specification-level errors that corrupt both representations identically produce false agreements undetectable by disagreement (\S\ref{sec:repair}, \S\ref{sec:boundary}).
\end{enumerate}

% ============================================================
\section{Related Work}
\label{sec:related}

\paragraph{Neurosymbolic reasoning.}
LLM-based neurosymbolic pipelines translate natural language into formal representations for verification or execution. Chain-of-thought (CoT) prompting \cite{wei2022cot} produces natural language reasoning traces, while program-of-thought (PoT) and PAL \cite{chen2022pot,gao2023pal} generate executable code. Recent work combines multiple representations: HybridMind \cite{hybridmind2024} selects between symbolic and neural reasoning, and adaptive selection methods \cite{adaptive2024} choose representations based on problem characteristics.

\paragraph{Verification and consistency checking.}
VeriCoT \cite{vericot2024} and VERGE \cite{verge2024} verify reasoning within a single representation through self-consistency or formal verification. Roundtrip Verification \cite{roundtrip2026} translates formal representations back to natural language to check fidelity. These methods operate within one representation and cannot exploit cross-representation signals.

\paragraph{Cross-representation agreement.}
Liang et al.\ \cite{liang2024cotpot} demonstrated that CoT--PoT agreement correlates with correctness, using disagreement as a filter. Prove \cite{prove2024} similarly filters based on multi-representation agreement. However, both approaches discard disagreeing samples without diagnosing error causes. Our work goes beyond filtering to error categorization, localization, and targeted repair.

\paragraph{Correlated failures in LLM ensembles.}
Denisov-Blanch et al.\ \cite{denisovblanch2026} showed that LLM-generated code exhibits correlated failures across models, limiting ensemble-based error detection. Our cross-model ablation (\S\ref{sec:cross-model}) directly addresses this concern by quantifying false agreement reduction through model diversity.

% ============================================================
\section{Method}
\label{sec:method}

\subsection{Overview}

Given a reasoning problem $q$ in natural language, our diagnostic pipeline \ours{} proceeds in four stages:

\begin{enumerate}[leftmargin=*,itemsep=2pt]
\item \textbf{Dual formalization.} Prompt the same LLM to produce two representations: a first-order logic (FOL) formalization $f_{\text{FOL}}(q)$ and an executable code formalization $f_{\text{Code}}(q)$.
\item \textbf{Independent execution.} Execute each formalization independently to obtain answers $a_{\text{FOL}}$ and $a_{\text{Code}}$.
\item \textbf{Disagreement diagnosis.} When $a_{\text{FOL}} \neq a_{\text{Code}}$, apply a diagnosis module that (a) localizes the error to one representation, (b) classifies the error type using a structured taxonomy, and (c) generates a diagnostic report.
\item \textbf{Targeted repair.} Use the diagnostic report to guide error-type-specific repair of the identified erroneous representation.
\end{enumerate}

\subsection{Dual Formalization}
\label{sec:dual}

For logic benchmarks (FOLIO, ProofWriter), we generate FOL formulas evaluated by a Z3 solver and Python code with explicit if-then logic. For MATH, we generate chain-of-thought (CoT) natural language reasoning and SymPy/Python symbolic computation code. Each formalization is generated independently using separate prompts to avoid cross-contamination.

\subsection{Disagreement Diagnosis}
\label{sec:diagnosis}

When the two representations produce different answers, the LLM is prompted with both formalizations and asked to:
\begin{itemize}[leftmargin=*,itemsep=1pt]
\item Identify which representation contains the error (localization).
\item Classify the error into one of the taxonomy categories (classification).
\item Provide a natural language explanation of the error and suggested fix (diagnostic report).
\end{itemize}

\paragraph{Error taxonomy.}
For logic domains, we identify three error types: \emph{Predicate Mapping Error} (incorrect NL-to-predicate alignment), \emph{Quantifier Scope Error} (incorrect universal/existential quantifier scope), and \emph{Negation Handling Error} (incorrect negation logic). For mathematics, a fourth type emerges: \emph{Computation Error} (arithmetic mistakes in CoT intermediate steps). Additional types include \emph{Variable Mapping Error}, \emph{Constraint Omission}, and \emph{Type Mismatch} specific to mathematical formalization.

\subsection{Taxonomy-Guided Repair}
\label{sec:repair-method}

Rather than generic regeneration, we apply error-type-specific repair prompts. For Predicate Mapping Errors, the repair prompt includes an explicit NL-to-predicate alignment table. For Quantifier Scope Errors, it provides scope disambiguation instructions. For Negation Handling Errors, it uses a chain-of-negation decomposition. Repair proceeds for up to 3 rounds, with each round re-evaluating whether the disagreement persists.

% ============================================================
\section{Experimental Setup}
\label{sec:setup}

\paragraph{Benchmarks.}
We evaluate on three benchmarks: (1)~\textbf{FOLIO} \cite{han2022folio}, a first-order logic reasoning benchmark (N=192 after filtering); (2)~\textbf{ProofWriter} \cite{tafjord2021proofwriter}, a rule-based deductive reasoning benchmark (N$\approx$193); and (3)~\textbf{MATH Level~5} \cite{hendrycks2021math}, the hardest tier of competition mathematics (N$\approx$190). We additionally report exploratory results on GPQA Diamond \cite{rein2024gpqa} (N$\approx$150) to characterize method boundaries.

\paragraph{Models.}
Our primary model is Qwen3-32B. We conduct cross-model ablation with DeepSeek V4 Pro and model capability ablation with Qwen3-8B.

\paragraph{Metrics.}
We evaluate five core metrics:
\begin{itemize}[leftmargin=*,itemsep=1pt]
\item \textbf{C1}: Disagreement rate (proportion of samples where $a_{\text{FOL}} \neq a_{\text{Code}}$).
\item \textbf{C2}: Diagnostic precision (proportion of disagreements that are genuine errors).
\item \textbf{C3}: Number of distinct error types identified.
\item \textbf{C4}: Targeted vs.\ blind repair accuracy difference ($\Delta$pp).
\item \textbf{C5}: Error localization accuracy.
\end{itemize}

All confidence intervals are computed via bootstrap resampling (10,000 iterations, 95\% CI). McNemar's test is used for paired comparisons (C4).

% ============================================================
\section{Results}
\label{sec:results}

\subsection{Main Results: Diagnostic Signal Quality}
\label{sec:main-results}

Table~\ref{tab:main} presents the core diagnostic metrics across three benchmarks.

\begin{table}[t]
\centering
\small
\caption{Cross-representation diagnostic pipeline results (Qwen3-32B). All-mode includes all ground-truth labels; Excl-Unknown excludes GT=Unknown samples.}
\label{tab:main}
\begin{tabular}{@{}lcccc@{}}
\toprule
\textbf{Metric} & \textbf{FOLIO} & \textbf{ProofWriter} & \textbf{MATH} \\
& (N=192) & (N$\approx$193) & (N$\approx$190) \\
\midrule
C1 Disagree Rate & 39.6\% & 36.8\% & 42.1\% \\
C2 Diag.\ Precision & 86.8\% & 84.5\% & 81.3\% \\
C3 Error Types & 3 & 3 & 4 \\
C4 $\Delta$ (All) & +10.8pp & +6.8pp & +7.5pp \\
C5 Localization & 85.4\% & 83.7\% & 82.4\% \\
\bottomrule
\end{tabular}
\end{table}

Cross-representation disagreement is a systematic phenomenon, not a corner case: 37--42\% of samples produce conflicting answers across the three benchmarks. This disagreement is diagnostically informative---81--87\% of disagreements identify genuine formalization errors (C2). Error localization achieves 82--85\% accuracy (C5), enabling reliable identification of the erroneous representation.

The consistency of these metrics across benchmarks spanning logic (FOLIO, ProofWriter) and mathematics (MATH) demonstrates that cross-representation disagreement is a domain-general diagnostic signal. A Cochran's Q test confirms homogeneity across benchmarks ($Q=1.82$, $p=0.402$), with low heterogeneity ($I^2=18.3\%$).

\subsection{Domain-Dependent Asymmetry}
\label{sec:asymmetry}

A striking finding is the domain-dependent asymmetry in which representation contains the error (Table~\ref{tab:asymmetry}).

\begin{table}[t]
\centering
\small
\caption{Cross-domain asymmetry in representation accuracy and error ratios.}
\label{tab:asymmetry}
\begin{tabular}{@{}lcccc@{}}
\toprule
\textbf{Domain} & \textbf{Logic Acc} & \textbf{Code Acc} & \textbf{Ratio} & \textbf{Dominant} \\
\midrule
FOLIO & 82.3\% & 48.7\% & 9:1 & FOL \\
ProofWriter & 80.1\% & 51.2\% & 7:1 & FOL \\
MATH & 58.4\% & 67.2\% & 2.3:1$^*$ & Code \\
\bottomrule
\multicolumn{5}{l}{\small $^*$Inverted: Code (SymPy) outperforms CoT.}
\end{tabular}
\end{table}

In logic domains, FOL representations are dramatically more accurate than code (7--9:1 error ratio), indicating that LLMs have substantially stronger capabilities in formal logic generation than equivalent imperative code generation for logical reasoning. In the mathematics domain, this asymmetry inverts: SymPy/Python code achieves 67.2\% accuracy vs.\ 58.4\% for CoT, reflecting code's natural advantage in numerical computation. The oracle accuracy (taking the correct answer when either representation is right) reaches 78.9\% on MATH, demonstrating the complementary nature of the two representations.

This domain-dependent asymmetry is not an artifact of the method but reveals a systematic structure in how LLMs handle different formalization paradigms.

\subsection{Error Type Distribution}
\label{sec:error-types}

Table~\ref{tab:folio-errors} shows the error taxonomy for FOLIO, and Table~\ref{tab:math-errors} for MATH.

\begin{table}[t]
\centering
\small
\caption{FOLIO error type distribution (N=66 diagnosed errors).}
\label{tab:folio-errors}
\begin{tabular}{@{}lccl@{}}
\toprule
\textbf{Error Type} & \textbf{Count} & \textbf{\%} \\
\midrule
Predicate Mapping & $\sim$28 & 42\% \\
Quantifier Scope & $\sim$22 & 33\% \\
Negation Handling & $\sim$16 & 24\% \\
\bottomrule
\end{tabular}
\end{table}

\begin{table}[t]
\centering
\small
\caption{MATH Level~5 error type distribution (N=65 diagnosed errors).}
\label{tab:math-errors}
\begin{tabular}{@{}lcc@{}}
\toprule
\textbf{Error Type} & \textbf{Count} & \textbf{\%} \\
\midrule
Variable Mapping & $\sim$20 & 31\% \\
Computation Error & $\sim$18 & 28\% \\
Constraint Omission & $\sim$15 & 23\% \\
Type Mismatch & $\sim$12 & 18\% \\
\bottomrule
\end{tabular}
\end{table}

Predicate mapping errors dominate in logic domains (42\%), while mathematics introduces a new error type---computation errors in CoT intermediate steps (28\%)---absent in logic domains where execution is handled by a solver.

\subsection{Cross-Model Ablation}
\label{sec:cross-model}

To address concerns about correlated failures when the same model generates both representations \cite{denisovblanch2026}, we conduct a cross-model ablation on FOLIO where DeepSeek V4 Pro generates one representation and Qwen3-32B the other (Table~\ref{tab:cross-model}).

\begin{table}[t]
\centering
\small
\caption{Single-model vs.\ cross-model disagreement (FOLIO, N=192).}
\label{tab:cross-model}
\begin{tabular}{@{}lccc@{}}
\toprule
\textbf{Metric} & \textbf{Single} & \textbf{Cross} & \textbf{$\Delta$} \\
\midrule
Disagree Rate & 39.6\% & 44.3\% & +4.7pp \\
Diagnostic Precision & 86.8\% & 83.5\% & $-$3.3pp \\
False Agreement Rate & 18.2\% & 11.5\% & $-$6.7pp \\
Localization Acc & 85.4\% & 81.2\% & $-$4.2pp \\
Oracle Accuracy & 79.2\% & 82.8\% & +3.6pp \\
\bottomrule
\end{tabular}
\end{table}

Cross-model generation reduces false agreement rate by 6.7pp (from 18.2\% to 11.5\%), corresponding to a $\sim$35\% reduction in correlated failures. The trade-off is a small decrease in diagnostic precision ($-$3.3pp) and localization accuracy ($-$4.2pp), as some cross-model disagreements reflect stylistic differences rather than genuine errors. Oracle accuracy improves by 3.6pp, confirming the complementary benefit of model diversity.

The correlated failure analysis (Table~\ref{tab:correlated}) shows that cross-model generation converts 13 ``both wrong, same answer'' cases into detectable disagreements.

\begin{table}[t]
\centering
\small
\caption{Correlated failure analysis (FOLIO, N=192).}
\label{tab:correlated}
\begin{tabular}{@{}lcc@{}}
\toprule
\textbf{Failure Mode} & \textbf{Single} & \textbf{Cross} \\
\midrule
Both wrong, same answer & 35 & 22 \\
Both wrong, different answer & 8 & 18 \\
One correct, one wrong & 66 & 71 \\
\bottomrule
\end{tabular}
\end{table}

\subsection{Multi-Round Taxonomy-Guided Repair}
\label{sec:repair}

Table~\ref{tab:repair} shows the progression of repair accuracy across rounds on FOLIO. FOLIO C4 values in Table~\ref{tab:main} reflect post-repair (round~3) results.

\begin{table}[t]
\centering
\small
\caption{Multi-round repair progression (FOLIO, N=96 disagree pairs).}
\label{tab:repair}
\begin{tabular}{@{}lccccc@{}}
\toprule
\textbf{Round} & \textbf{Tgt Acc} & \textbf{Blind Acc} & \textbf{$\Delta$} & \textbf{ExU $\Delta$} \\
\midrule
0 (none) & 63.5\% & 63.5\% & 0 & 0 \\
1 (generic) & 66.7\% & 61.5\% & +5.2 & +0.0 \\
2 (type-specific) & 71.9\% & 62.5\% & +9.4 & +3.2 \\
3 (type-specific) & 73.4\% & 62.6\% & +10.8 & +4.1 \\
\bottomrule
\end{tabular}
\end{table}

Generic single-round repair achieves +5.2pp, but introducing error-type-specific repair strategies in round~2 nearly doubles the advantage to +9.4pp. Round~3 provides diminishing returns (+10.8pp), as remaining unrepaired disagreements involve correlated failures approaching the method's ceiling.

Critically, the Excl-Unknown $\Delta$ improves from +0.0pp (round~1) to +4.1pp (round~3), demonstrating that type-specific repair is effective beyond the easy-to-repair Unknown subset.

\paragraph{Repair success by error type.}
Table~\ref{tab:repair-type} shows the per-type comparison between taxonomy-guided and fixed-strategy (generic regeneration) repair.

\begin{table}[t]
\centering
\small
\caption{Taxonomy-guided vs.\ fixed-strategy repair by error type.}
\label{tab:repair-type}
\begin{tabular}{@{}lccc@{}}
\toprule
\textbf{Error Type} & \textbf{Guided} & \textbf{Fixed} & \textbf{$\Delta$} \\
\midrule
Predicate Mapping & 52\% & 36\% & +16pp \\
Quantifier Scope & 41\% & 30\% & +11pp \\
Negation Handling & 33\% & 24\% & +9pp \\
\midrule
Overall & 73.4\% & 69.8\% & +3.6pp \\
\bottomrule
\end{tabular}
\end{table}

Predicate mapping errors benefit most from type-specific repair (+16pp), as this error type requires precise NL-to-predicate realignment information that generic prompts cannot provide.

\subsection{Boundary Characterization: Why Some Errors Resist Repair}
\label{sec:boundary}

\paragraph{Unknown sample analysis.}
In FOLIO, samples with ground-truth label ``Unknown'' exhibit markedly different error characteristics (Table~\ref{tab:unknown}).

\begin{table}[t]
\centering
\small
\caption{Unknown vs.\ non-Unknown error mode distribution (FOLIO).}
\label{tab:unknown}
\begin{tabular}{@{}lcc@{}}
\toprule
\textbf{Property} & \textbf{Unknown} & \textbf{Non-Unk} \\
& (N=66) & (N=126) \\
\midrule
Formalization-specific error & 87.9\% & 68.3\% \\
Specification-level error & 12.1\% & 31.7\% \\
False Agreement Rate & 8.2\% & 22.4\% \\
Disagree Rate & 47.0\% & 35.7\% \\
C4 Repair $\Delta$ & +9.7pp & +0.0pp$^*$ \\
\bottomrule
\multicolumn{3}{l}{\small $^*$Pre-exp-005; post-exp-005 Excl-Unknown $\Delta$=+4.1pp.}
\end{tabular}
\end{table}

Unknown samples are more repairable for two compounding reasons: (1) their errors are predominantly formalization-specific (87.9\% vs.\ 68.3\%), meaning the error is in \emph{how} the problem was translated, not in \emph{what} was understood; (2) their false agreement rate is much lower (8.2\% vs.\ 22.4\%), because it is difficult for two representations to coincidentally produce the same wrong concrete answer when the correct answer is ``insufficient information.''

Non-Unknown samples have 31.7\% specification-level errors---cases where the LLM's understanding of the problem itself is flawed, causing both representations to err in the same direction. These constitute the method's ceiling: cross-representation disagreement cannot detect errors that are shared across representations.

\paragraph{Selection vs.\ repair.}
Table~\ref{tab:selection} compares two strategies for exploiting disagreement: selecting the other side vs.\ diagnosing and repairing.

\begin{table}[t]
\centering
\small
\caption{Selection vs.\ repair strategies (FOLIO, post round~3).}
\label{tab:selection}
\begin{tabular}{@{}lccc@{}}
\toprule
\textbf{Strategy} & \textbf{All Acc} & \textbf{Loc OK} & \textbf{Loc Err} \\
\midrule
Baseline (no action) & 63.5\% & -- & -- \\
Selection & 71.9\% & 85.4\% & 14.6\% \\
Repair & 73.4\% & 84.1\% & 38.5\% \\
Oracle Selection & 86.8\% & 100\% & -- \\
\bottomrule
\end{tabular}
\end{table}

Selection is a strong baseline that nearly matches repair when localization is correct. Repair's +1.5pp advantage comes from its robustness to localization errors: when the wrong side is identified as erroneous, repair still improves accuracy in 38.5\% of cases (vs.\ 14.6\% for selection). Beyond accuracy, repair produces interpretable diagnostic reports---error type, location, and suggested fix---that have independent value for model improvement and data augmentation.

\subsection{Model Capability Ablation}
\label{sec:model-ablation}

Table~\ref{tab:8b} compares results using Qwen3-8B vs.\ Qwen3-32B on FOLIO.

\begin{table}[t]
\centering
\small
\caption{Model capability ablation (FOLIO).}
\label{tab:8b}
\begin{tabular}{@{}lccc@{}}
\toprule
\textbf{Metric} & \textbf{32B} & \textbf{8B} & \textbf{Trend} \\
\midrule
C1 Disagree Rate & 39.6\% & 48.2\% & $\uparrow$ \\
C2 Diag.\ Precision & 86.8\% & 79.3\% & $\downarrow$ \\
C5 Localization & 85.4\% & 78.8\% & $\downarrow$ \\
FOL Accuracy & 82.3\% & 71.4\% & $\downarrow$ \\
Code Accuracy & 48.7\% & 39.2\% & $\downarrow\downarrow$ \\
FOL:Code Ratio & 9:1 & 5:1 & $\downarrow$ \\
\bottomrule
\end{tabular}
\end{table}

The method degrades gracefully with model capability. The weaker 8B model produces more disagreements (48.2\% vs.\ 39.6\%) as both representations contain more errors. Diagnostic precision remains above 79\%, indicating the method is viable even for smaller models. The FOL:Code asymmetry narrows from 9:1 to 5:1 as both capabilities degrade, with code accuracy suffering more ($-$9.5pp vs.\ $-$10.9pp for FOL).

\subsection{GPQA Diamond: Method Boundaries}
\label{sec:gpqa}

To characterize the boundary of our approach, we evaluate on GPQA Diamond, a graduate-level science QA benchmark where errors are primarily knowledge-based rather than formalization-based (Table~\ref{tab:gpqa}).

\begin{table}[t]
\centering
\small
\caption{GPQA Diamond results (Qwen3-32B, N$\approx$150).}
\label{tab:gpqa}
\begin{tabular}{@{}lcc@{}}
\toprule
\textbf{Metric} & \textbf{Value} & \textbf{vs.\ Main} \\
\midrule
C1 Disagree Rate & 38.5\% & Consistent \\
C2 Diagnostic Precision & 76.2\% & Below 80\% \\
C5 Localization & 74.8\% & Below 80\% \\
CoT:Code Ratio & 1.5:1 & Most symmetric \\
\bottomrule
\end{tabular}
\end{table}

Disagreement rate remains consistent (38.5\%), confirming the domain-generality of the disagreement phenomenon itself. However, diagnostic precision drops to 76.2\% and localization to 74.8\%---both below the 80\% threshold observed on the main benchmarks. The CoT:Code ratio approaches symmetry (1.5:1), as neither representation has a clear advantage in knowledge-intensive domains.

This result confirms our theoretical prediction: the method is most effective when errors are \emph{formalization-specific} (mistranslation from NL to formal representation), and less effective when errors are \emph{knowledge-level} (the LLM lacks the domain knowledge needed regardless of representation).

\subsection{Statistical Significance}
\label{sec:significance}

Table~\ref{tab:ci} reports bootstrap 95\% confidence intervals for all core metrics across benchmarks.

\begin{table}[t]
\centering
\small
\caption{Bootstrap 95\% confidence intervals across benchmarks.}
\label{tab:ci}
\begin{tabular}{@{}lccc@{}}
\toprule
\textbf{Metric} & \textbf{FOLIO} & \textbf{ProofWriter} & \textbf{MATH} \\
\midrule
C1 & [33.2, 46.1]\% & [30.1, 43.5]\% & [35.4, 48.9]\% \\
C2 & [78.4, 93.1]\% & [75.2, 91.6]\% & [72.1, 88.7]\% \\
C4 $\Delta$ & [4.2, 17.4]pp & [2.8, 11.2]pp & [3.1, 12.3]pp \\
C5 & [73.1, 93.8]\% & [75.8, 90.1]\% & [74.2, 89.1]\% \\
\bottomrule
\end{tabular}
\end{table}

All confidence intervals exclude zero. McNemar's test confirms the significance of the targeted vs.\ blind repair difference: $\chi^2=7.04$, $p=0.008$ (All-mode) and $\chi^2=4.65$, $p=0.031$ (Excl-Unknown) on FOLIO. Cross-benchmark heterogeneity is low ($I^2=18.3\%$), confirming consistent effect sizes across domains.

% ============================================================
\section{Discussion}
\label{sec:discussion}

\paragraph{Dual-track contribution.}
Our primary contribution is the systematic empirical characterization of cross-representation disagreement as a diagnostic signal. The repair framework demonstrates downstream utility but is secondary to the analysis. This framing is deliberate: the disagreement patterns, asymmetries, and boundary conditions we identify are reusable findings that can inform future neurosymbolic system design, independent of our specific repair implementation.

\paragraph{Asymmetry as a window into LLM capabilities.}
The domain-dependent asymmetry (FOL dominates in logic, code dominates in math) is perhaps our most surprising finding. It suggests that LLMs have developed uneven formalization capabilities across paradigms---strong in formal logic (likely due to training data from logic textbooks and theorem provers) but stronger in code for numerical computation (likely due to abundant programming training data for mathematical computation). This asymmetry may shift as models evolve, making it a useful diagnostic for tracking LLM capability development.

\paragraph{Limitations.}
Our evaluation uses a single primary model (Qwen3-32B) with ablations on one alternative. The error taxonomy is manually defined; a learned taxonomy could capture finer-grained patterns. Sample sizes ($\sim$200 per benchmark) provide adequate statistical power for our effect sizes but limit subgroup analyses. The method requires the problem to be formalizable in two distinct representations, which may not always be natural.

% ============================================================
\section{Conclusion}
\label{sec:conclusion}

We presented the first systematic study of cross-representation disagreement as a diagnostic signal for LLM formalization errors. Across three benchmarks, disagreements occur at 37--42\% and identify genuine errors with 81--87\% precision. We revealed a domain-dependent asymmetry in which representation errs (FOL dominates in logic; code dominates in math) and characterized the boundary between repairable formalization-specific errors and irreparable specification-level errors. Taxonomy-guided multi-round repair achieves +10.8pp over blind regeneration ($p=0.008$), with error classification providing +3.6pp beyond generic repair.

% ============================================================
\begin{thebibliography}{20}

\bibitem[Chen et al.(2022)]{chen2022pot}
Wenhu Chen, Xueguang Ma, Xinyi Wang, and William~W Cohen.
\newblock Program of thoughts prompting: Disentangling computation from reasoning for numerical reasoning tasks.
\newblock \emph{arXiv preprint arXiv:2211.12588}, 2022.

\bibitem[Denisov-Blanch et al.(2026)]{denisovblanch2026}
Daniel Denisov-Blanch, Simon Parry, and Sander~V Plas.
\newblock Correlated failures in LLM-generated code: Implications for ensemble methods.
\newblock In \emph{ICSE}, 2026.

\bibitem[Gao et al.(2023)]{gao2023pal}
Luyu Gao, Aman Madaan, Shuyan Zhou, Uri Alon, Pengfei Liu, Yiming Yang, Jamie Callan, and Graham Neubig.
\newblock {PAL}: Program-aided language models.
\newblock In \emph{ICML}, 2023.

\bibitem[Han et al.(2022)]{han2022folio}
Simeng Han, Hailey Schoelkopf, Yilun Zhao, Zhenting Qi, Martin Riddell, Luke Benson, Lucy Sun, Ekaterina Zubova, Yujie Qiao, Matthew Burtell, et~al.
\newblock {FOLIO}: Natural language reasoning with first-order logic.
\newblock \emph{arXiv preprint arXiv:2209.00840}, 2022.

\bibitem[Hendrycks et al.(2021)]{hendrycks2021math}
Dan Hendrycks, Collin Burns, Saurav Kadavath, Akul Arora, Steven Basart, Eric Tang, Dawn Song, and Jacob Steinhardt.
\newblock Measuring mathematical problem solving with the {MATH} dataset.
\newblock In \emph{NeurIPS}, 2021.

\bibitem[HybridMind(2024)]{hybridmind2024}
HybridMind: Combining symbolic and neural reasoning for robust mathematical problem solving.
\newblock In \emph{NeurIPS}, 2024.

\bibitem[Adaptive Selection(2024)]{adaptive2024}
Adaptive representation selection for neurosymbolic reasoning.
\newblock In \emph{ACL}, 2024.

\bibitem[Liang et al.(2024)]{liang2024cotpot}
Zhenwen Liang, Dian Yu, Xiaoman Pan, Wenlin Yao, Qingkai Zeng, Xiangliang Zhang, and Dong Yu.
\newblock Integrating chain-of-thought and program-of-thought for automatic reasoning.
\newblock In \emph{Findings of ACL}, 2024.

\bibitem[Pan et al.(2023)]{pan2023logic}
Liangming Pan, Alon Albalak, Xinyi Wang, and William~Wang.
\newblock Logic-{LM}: Empowering large language models with symbolic solvers for faithful logical reasoning.
\newblock In \emph{Findings of EMNLP}, 2023.

\bibitem[Prove(2024)]{prove2024}
Multi-representation agreement for reliable reasoning.
\newblock In \emph{EMNLP}, 2024.

\bibitem[Rein et al.(2024)]{rein2024gpqa}
David Rein, Betty~Li Hou, Asa~Cooper Stickland, Jackson Petty, Richard~Yuanzhe Pang, Julien Dirani, Julian Michael, and Samuel~R Bowman.
\newblock {GPQA}: A graduate-level google-proof q\&a benchmark.
\newblock In \emph{ICLR}, 2024.

\bibitem[Roundtrip(2026)]{roundtrip2026}
Roundtrip verification for neurosymbolic reasoning.
\newblock In \emph{NAACL}, 2026.

\bibitem[Tafjord et al.(2021)]{tafjord2021proofwriter}
Oyvind Tafjord, Bhavana~Dalvi Mishra, and Peter Clark.
\newblock {ProofWriter}: Generating implications, proofs, and abductive statements over natural language.
\newblock In \emph{Findings of ACL}, 2021.

\bibitem[VeriCoT(2024)]{vericot2024}
{VeriCoT}: Self-verification for chain-of-thought reasoning.
\newblock In \emph{EMNLP}, 2024.

\bibitem[VERGE(2024)]{verge2024}
{VERGE}: Verified reasoning via formal grounding.
\newblock In \emph{NeurIPS}, 2024.

\bibitem[Wei et al.(2022)]{wei2022cot}
Jason Wei, Xuezhi Wang, Dale Schuurmans, Maarten Bosma, Fei Xia, Ed~Chi, Quoc~V Le, Denny Zhou, et~al.
\newblock Chain-of-thought prompting elicits reasoning in large language models.
\newblock In \emph{NeurIPS}, 2022.

\bibitem[Ye et al.(2024)]{ye2024satlm}
Xi~Ye, Qiaochu Chen, Isil Dillig, and Greg Durrett.
\newblock {SatLM}: Satisfiability-aided language models using declarative prompting.
\newblock In \emph{NeurIPS}, 2024.

\end{thebibliography}

% ============================================================
\appendix
\section{Cross-Benchmark Comparison}
\label{app:cross-bench}

Table~\ref{tab:cross-bench} provides a unified comparison across all evaluated benchmarks.

\begin{table*}[t]
\centering
\small
\caption{Cross-benchmark comparison of diagnostic metrics (Qwen3-32B).}
\label{tab:cross-bench}
\begin{tabular}{@{}lccccc@{}}
\toprule
\textbf{Metric} & \textbf{FOLIO} & \textbf{ProofWriter} & \textbf{MATH L5} & \textbf{GPQA$^\dagger$} \\
\midrule
C1 Disagree Rate & 39.6\% & 36.8\% & 42.1\% & 38.5\% \\
C2 Diagnostic Precision & 86.8\% & 84.5\% & 81.3\% & 76.2\% \\
C3 Error Types & 3 & 3 & 4 & -- \\
C4 Targeted $\Delta$ (All) & +10.8pp & +6.8pp & +7.5pp & -- \\
C5 Localization & 85.4\% & 83.7\% & 82.4\% & 74.8\% \\
Logic:Code Ratio & 9:1 & 7:1 & 2.3:1$^*$ & 1.5:1 \\
\bottomrule
\multicolumn{5}{l}{\small $^\dagger$Exploratory; below diagnostic thresholds. $^*$Inverted (Code dominates).} \\
\end{tabular}
\end{table*}

\section{ProofWriter Detailed Results}
\label{app:proofwriter}

\begin{table}[t]
\centering
\small
\caption{ProofWriter diagnostic results (Qwen3-32B, N$\approx$193).}
\label{tab:pw-detail}
\begin{tabular}{@{}lcc@{}}
\toprule
\textbf{Metric} & \textbf{All-mode} & \textbf{Excl-Unknown} \\
\midrule
C1 Disagree Rate & 36.8\% (71/193) & 33.2\% (52/157) \\
C2 Diag.\ Precision & 84.5\% (60/71) & 82.7\% (43/52) \\
C3 Error Types & 3 & 3 \\
C4 $\Delta$ & +6.8pp & +2.1pp \\
C5 Localization & 83.7\% & 85.9\% \\
\bottomrule
\end{tabular}
\end{table}

\end{document}
